\chapter{Backtests événementiels}\label{ch:rs:event-driven-backtests}

Une stratégie qui laisse reposer des ordres à cours limité à un tick de part et d'autre du dernier cours voit ses ordres exécutés, dans un \omterm{def:rs:vectorised-backtests:backtest}{backtest} sur barres, chaque fois que l'amplitude d'une barre atteint
leur prix. Si l'on exige à la place que le cours traverse l'ordre d'un tick avant de compter une exécution, le même \omterm{def:rs:vectorised-backtests:backtest}{backtest} n'exécute plus que 55~\% aussi souvent et transforme
un petit profit quotidien en une perte de 584~\$ par jour, les exécutions étant suivies d'un mouvement de deux ticks contre elles. Il suffit d'ajouter une latence de saisie des ordres de cinq secondes,
et la réponse du \omterm{def:rs:vectorised-backtests:backtest}{backtest} dépend d'une question à laquelle les barres ne peuvent pas répondre~: le cours a-t-il touché l'ordre pendant qu'il était actif~? Le niveau~1 du
chapitre~16 ne permet même pas d'exprimer une telle stratégie. Le \omterm{def:rs:event-driven-backtests:evbt}{backtest événementiel}, lui, le permet~: il simule les ordres sur une horloge, et il fait de chaque hypothèse
sur les exécutions un modèle explicite et remplaçable. Ce chapitre le construit (\texttt{firm.evbt}) et montre quelle part du résultat d'une stratégie passive revient au
\omterm{def:rs:event-driven-backtests:fill}{modèle d'exécution}.

\section{Architecture~: horloge, événements, gestionnaires}

\begin{definition}[Backtest événementiel, horloge simulée]\label{def:rs:event-driven-backtests:evbt}
Un \emph{backtest événementiel}\index{backtest événementiel} traite un flux d'événements ordonnés dans le temps (données de marché, accusés de réception des ordres,
exécutions, annulations, opérations sur titres) un à un, chacun modifiant l'état de la stratégie, de ses ordres et de son portefeuille. Son
\emph{horloge simulée}\index{horloge simulée} est l'instant de l'événement en cours de traitement~; aucune information horodatée plus tard ne doit être utilisée.
\end{definition}

\begin{omfigure}
\begin{tikzpicture}[font=\small,
  box/.style={draw=omDef,fill=omDef!8,rounded corners=2pt,align=center,inner sep=4pt,minimum height=9mm,text width=2.7cm},
  arr/.style={-Stealth,thick,draw=black!60}]
\node[box] (q) at (0,0) {file d'événements\\ (instant, type, numéro de séquence)};
\node[box] (d) at (3.6,0) {répartiteur~:\\ fait avancer l'horloge};
\node[box] (m) at (7.2,0) {gestionnaire des ordres~:\\ cycle de vie};
\node[box] (p) at (10.8,0) {portefeuille~:\\ position, trésorerie};
\node[box] (s) at (7.2,2.0) {stratégie~:\\ on\_bar, on\_fill};
\node[box] (f) at (7.2,-2.0) {modèle d'exécution~: contact,\\ pénétration, plafond};
\draw[arr] (q) -- (d);
\draw[arr] (d) -- (m);
\draw[arr] (m) -- (p);
\draw[arr] (d) |- (s);
\draw[arr] (s) -- node[right,font=\footnotesize]{soumission, annulation} (m);
\draw[arr] (m) -- (f);
\draw[arr] (f) -| (p);
\draw[arr,dashed] (s.north) -- ++(0,0.5) -| node[pos=0.25,above,font=\footnotesize]{accusés de réception et annulations, après la latence} (q.north);
\end{tikzpicture}

\omcaption{Architecture de \texttt{firm.evbt}~: une file d'événements ordonnée par instant, type et numéro de séquence, un répartiteur et des gestionnaires qui
ne peuvent programmer des événements que dans le futur.}\label{fig:rs:event-driven-backtests:architecture}
\end{omfigure}

\texttt{firm.evbt} gère une seule file de priorité. À la fin de chaque barre, le répartiteur confronte d'abord les ordres actifs à la barre par le
\omterm{def:rs:event-driven-backtests:fill}{modèle d'exécution}, puis valorise le portefeuille, puis appelle la stratégie, qui peut soumettre ou annuler des ordres~; une soumission ne devient un ordre actif qu'au moment où son
accusé de réception arrive, après la latence de saisie des ordres, et une annulation ne prend effet que lorsque son propre événement arrive. Les événements au même
instant sont classés par type (barres avant événements d'ordres), puis par numéro de séquence, de sorte qu'un passage du \omterm{def:rs:vectorised-backtests:backtest}{backtest} soit déterministe~: mêmes données, stratégie et
modèles donnent les mêmes exécutions au dernier chiffre près. Le résultat est le \texttt{BacktestResult} du chapitre~16, avec les ordres et les exécutions qui y sont joints.

\section{Ordres et exécutions à l'échelle de la barre}

\begin{definition}[Cycle de vie de l'ordre, exécution partielle]\label{def:rs:event-driven-backtests:lifecycle}
Le \emph{cycle de vie de l'ordre}\index{cycle de vie de l'ordre} est la suite des états par lesquels passe un ordre~: en attente (envoyé, pas encore accusé de réception),
actif, partiellement exécuté, puis exécuté ou annulé. Une \emph{exécution partielle}\index{exécution partielle} exécute une partie de la quantité d'un ordre, laissant le
reste actif.
\end{definition}

\begin{definition}[Modèle d'exécution, exécution au contact, exécution par pénétration, plafond de participation au volume]\label{def:rs:event-driven-backtests:fill}
Un \emph{modèle d'exécution}\index{modèle d'exécution} décide, à partir des données de marché, si un ordre actif s'exécute, à quel prix et pour quelle quantité. La
\emph{exécution au contact}\index{exécution au contact} exécute un ordre à cours limité en totalité lorsqu'une barre échange à son prix~; la \emph{exécution par pénétration}\index{exécution par pénétration}
seulement lorsque la barre échange un nombre donné de ticks au-delà de ce prix~; un \emph{plafond de participation au volume}\index{plafond de participation au volume} limite la quantité
exécutée dans une barre à une fraction du volume de la barre.
\end{definition}

Une barre enregistre quatre prix et un volume~; elle n'enregistre pas où se trouvait l'ordre de la stratégie dans la file à son prix, combien a été négocié à ce prix, ni
si les transactions à ce prix ont eu lieu avant ou après l'arrivée de l'ordre. Chaque \omterm{def:rs:event-driven-backtests:fill}{modèle d'exécution} comble cette lacune par une hypothèse. L'\omterm{def:rs:event-driven-backtests:fill}{exécution au contact} suppose
que l'ordre était en tête de file. L'\omterm{def:rs:event-driven-backtests:fill}{exécution par pénétration} suppose qu'il était en queue de file~: si le cours a traversé ce prix, tout ce qui le précédait a forcément été négocié.
Le \omterm{def:rs:event-driven-backtests:fill}{plafond de participation au volume} suppose que l'ordre peut prendre une part fixe du volume de la barre~; le moteur de \omterm{def:rs:vectorised-backtests:backtest}{backtest} open source Zipline, avec son modèle \texttt{VolumeShareSlippage}
modèle, par exemple, plafonne les exécutions à 2.5~\% du volume de chaque barre par défaut et les valorise à partir du cours de clôture de la barre avec un impact quadratique. La vérité se situe
entre les \omterm{def:rs:event-driven-backtests:fill}{exécutions au contact} et par pénétration, et dépend de la file (chapitre~18).

La stratégie du chapitre (\texttt{rs\_evbt.Quoter}) tourne sur des barres d'une minute construites à partir de quatre journées simulées de 6.5~heures de \texttt{firm.tape}~: à chaque
clôture de barre, elle annule ses ordres et laisse reposer un lot (100 actions) à un tick sous le cours de clôture et un lot à un tick au-dessus, dans la limite d'un inventaire de cinq
lots, en valorisant au cours de clôture. Par jour, en moyenne~:

\begin{center}
\small
\begin{tabular}{@{}p{6.0cm}cccc@{}}
\toprule
\omterm{def:rs:event-driven-backtests:fill}{modèle d'exécution} (sans latence) & lots exécutés & taux d'exécution & P\&L (\$) & mark-out (ticks)\\
\midrule
\omterm{def:rs:event-driven-backtests:fill}{exécution au contact} & 523 & 70~\% & $+17$ & $-0.16$\\
\omterm{def:rs:event-driven-backtests:fill}{exécution par pénétration} d'un tick & 290 & 41~\% & $-584$ & $-2.00$\\
\omterm{def:rs:event-driven-backtests:fill}{exécution au contact}, plafonnée à 2.5~\% du volume & 388 & 55~\% & $+20$ & $-0.18$\\
niveau 1~: achat après une barre baissière, vente après une barre haussière & --- & --- & $-215$ & ---\\
\bottomrule
\end{tabular}
\end{center}

Le mark-out (le cours de clôture cinq barres après une exécution, dans le sens de l'exécution) mesure la sélection adverse~: les \omterm{def:rs:event-driven-backtests:fill}{exécutions par pénétration} sont celles qui surviennent
lorsque le cours traverse l'ordre, et elles perdent en moyenne deux ticks en cinq minutes (voir la \cref{fig:rs:event-driven-backtests:models}). Le
P\&L quotidien de l'\omterm{def:rs:event-driven-backtests:fill}{exécution au contact} (de $-\$333$ à $+\$206$ sur les quatre jours) reste dans son bruit autour de zéro~; celui de l'\omterm{def:rs:event-driven-backtests:fill}{exécution par pénétration} est négatif chaque jour. Un taux
d'exécution de 70~\% pour un ordre situé à un tick du cours est en soi un signal d'alarme~: un ordre réel rejoint la fin d'une file qui peut compter des milliers d'actions.

\section{Instant de décision et latence}

\begin{definition}[Latence et politique de l'ordre actif dans la barre]\label{def:rs:event-driven-backtests:latency}
Avec une latence de saisie des ordres, un ordre envoyé à la clôture d'une barre est actif à partir d'un instant quelconque de la barre suivante. Une politique \emph{conservatrice} laisse un ordre
s'exécuter dans une barre que s'il a été actif pendant toute la barre~; une politique \emph{optimiste}, s'il l'a été à un moment quelconque de celle-ci. Une barre ne peut dire laquelle des deux est juste,
car elle ne dit pas à quel moment de la barre le cours a atteint l'ordre.
\end{definition}

Avec une latence de cinq secondes, les ordres du quoter, annulés et remplacés à chaque clôture de barre, ne sont jamais actifs pendant toute une barre~: sous la
politique conservatrice, ils ne s'exécutent jamais. Sous la politique optimiste, ils s'exécutent plus souvent que sans latence (585 lots par jour pour l'\omterm{def:rs:event-driven-backtests:fill}{exécution au contact}), et
perdent davantage~: $-\$41$ par jour pour l'\omterm{def:rs:event-driven-backtests:fill}{exécution au contact}, $-\$776$ pour l'\omterm{def:rs:event-driven-backtests:fill}{exécution par pénétration}, $-\$190$ pour l'exécution plafonnée, le mark-out des exécutions plafonnées passant
de $-0.18$ à $-0.87$ ticks. Les deux politiques encadrent la réponse, et une stratégie dont le résultat passe de rien à une perte entre les deux ne peut être
jugée au niveau~2. C'est le signal qu'il faut passer au niveau~3, où l'arrivée de l'ordre, sa position dans la file et chaque transaction à son prix sont des événements
à part entière.

\begin{omfigure}
\begin{tikzpicture}
\begin{axis}[name=lots,width=7cm,height=5.4cm,ybar,bar width=9pt,title={\small lots exécutés par jour},tick label style={font=\footnotesize},
  xtick={0,1,2,3,4,5},xticklabels={contact,pénétr.,plafonné,contact,pénétr.,plafonné},x tick label style={rotate=45,anchor=east},
  xmin=-0.6,xmax=5.6,ymin=0,ymax=760,grid=major,grid style={gray!25},table/col sep=comma]
\addplot[fill=omDef!55,draw=omDef] table[x=k,y=lots]{figdata/research/17-event-driven-backtests/models.csv};
\draw[black!50,dashed] (axis cs:2.5,0) -- (axis cs:2.5,760);
\node[font=\footnotesize,anchor=north] at (axis cs:1,750) {sans latence};
\node[font=\footnotesize,anchor=north] at (axis cs:4,750) {5 s, optimiste};
\end{axis}
\begin{axis}[at={(lots.east)},anchor=west,xshift=1.4cm,width=7cm,height=5.4cm,ybar,bar width=9pt,title={\small P\&L par jour (\$)},
  tick label style={font=\footnotesize},xtick={0,1,2,3,4,5,6},xticklabels={contact,pénétr.,plafonné,contact,pénétr.,plafonné,niveau 1},
  x tick label style={rotate=45,anchor=east},ymin=-850,ymax=100,grid=major,grid style={gray!25},table/col sep=comma]
\addplot[fill=omThm!45,draw=omThm] table[x=k,y=pnl]{figdata/research/17-event-driven-backtests/models.csv};
\draw[black!50,dashed] (axis cs:2.5,-850) -- (axis cs:2.5,100);
\end{axis}
\end{tikzpicture}

\omcaption{Le teneur de cotation passif sur quatre journées simulées sous chaque \omterm{def:rs:event-driven-backtests:fill}{modèle d'exécution}, sans latence et avec une latence de cinq secondes sous la politique optimiste
(sous la politique prudente, il n'est jamais exécuté), ainsi que la version au niveau 1 de la même idée. Données~: \texttt{rs\_evbt} sur les barres de \texttt{firm.tape}.}\label{fig:rs:event-driven-backtests:models}
\end{omfigure}

\section{Calendriers et opérations sur titres dans la boucle}

Les événements qu'un \omterm{def:rs:vectorised-backtests:backtest}{backtest vectorisé} intègre dans les cours ajustés sont, dans un \omterm{def:rs:event-driven-backtests:evbt}{backtest événementiel}, des événements placés dans la file. Un split multiplie la position et les
quantités des ordres en attente par son ratio et divise leurs cours limites par ce même ratio, au moment où il prend effet~; les barres conservent les cours tels qu'ils ont été échangés, de sorte que
un ordre à cours limité qui a survécu à un split non ajusté se retrouverait au double du nouveau cours et serait exécuté instantanément. Un dividende verse un montant en espèces à la date de paiement aux positions
détenues à la date d'enregistrement. Les demi-journées mettent fin à la séance plus tôt, les enchères remplacent la négociation en continu à l'ouverture et à la clôture (livre~1, chapitre~13), et une suspension
gèle le carnet d'ordres~; chacun doit figurer dans le calendrier que lit le répartiteur, avec son propre gestionnaire, et chacun a déjà conduit un \omterm{def:rs:vectorised-backtests:backtest}{backtest} à trader à un moment où le marché
était fermé. \texttt{firm.evbt} implémente les splits et laisse les autres pour la suite de la construction~; son test d'acceptation détient un lot à travers un split de
deux pour un et vérifie que la position double et que le capital ne bouge pas.

\section{Déterminisme}

Un \omterm{def:rs:vectorised-backtests:backtest}{backtest} qui donne des réponses différentes d'une exécution à l'autre ne peut être ni débogué, ni comparé, ni relu. Le déterminisme exige trois choses~: un ordre total sur les événements
(instant, puis nature, puis numéro de séquence~; jamais l'ordre dans lequel un dictionnaire se trouve itéré par hasard), aucun aléa sauf celui de générateurs initialisés par une graine passés
explicitement, et aucune dépendance à l'horloge murale. Les tests de \texttt{firm.evbt} exécutent le même \omterm{def:rs:vectorised-backtests:backtest}{backtest} deux fois et comparent les rendements nets bit à bit. La
même discipline permet de comparer un \omterm{def:rs:vectorised-backtests:backtest}{backtest} de niveau 2 et un \omterm{def:rs:vectorised-backtests:backtest}{backtest} de niveau 3 sur les mêmes événements, et de rejouer une exécution en direct à partir de ses journaux (chapitre~19).

\section{Tutoriel~: l'ordre à cours limité qui était toujours exécuté}\label{tut:rs:event-driven-backtests:tut}

\begin{tutorial}
\textbf{Objectif.} Lancer le teneur de cotation passif dans \texttt{firm.evbt} sous trois \omterm{def:rs:event-driven-backtests:fill}{modèles d'exécution} et deux latences, avec les deux politiques de présence dans la barre, et comparer avec
la version au niveau 1 de la même idée. \textbf{État final~:} la \cref{fig:rs:event-driven-backtests:models}~; le tableau.
\begin{tutsteps}
\item \textbf{Les \omterm{def:rs:event-driven-backtests:fill}{modèles d'exécution}}~: au contact, par pénétration, et un plafond sur le volume de la barre.
\omcode{code/firm/evbt/firm_evbt.py}{70}{104}{Trois modèles d'exécution pour les barres.}\label{lst:rs:event-driven-backtests:fills}
\item \textbf{Appariement}~: quels ordres étaient actifs dans la barre, et ce que le modèle exécute.
\omcode{code/firm/evbt/firm_evbt.py}{151}{175}{Appariement des ordres en attente avec une barre.}\label{lst:rs:event-driven-backtests:match}
\item \textbf{Exécuter} \texttt{summary()}, \texttt{level1()} et \texttt{fig\_evbt.py}.
\end{tutsteps}
\textbf{Pour aller plus loin.} Placer les ordres à la clôture plutôt qu'à un tick de distance et observer le taux d'exécution du modèle au contact se rapprocher de 100~\%~; ajouter une
règle d'annulation seulement si le cours a bougé, afin que les ordres survivent entiers aux barres sous latence.
\end{tutorial}

\section{Construction~: le backtester de niveau 2}\label{bld:rs:event-driven-backtests:evbt}

\begin{build}
\bfield{Rôle} Le backtester de la société pour les stratégies dont les ordres, les exécutions et le calendrier comptent à la résolution de la barre~: ordres à cours limité, ordres stop, calendriers
intrajournaliers, opérations sur titres dans la boucle.
\bfield{Interface} \texttt{Order}, \texttt{Fill}, \texttt{TouchFill()}, \texttt{PenetrationFill(ticks)}, \texttt{VolumeCapFill(inner, participation)},
\texttt{Strategy.on\_bar(ctx, i, bar)}, \texttt{Strategy.on\_fill(ctx, fill)},
\texttt{Engine(bars, strategy, fill\_model, latency, fee, capital, splits, policy).run()} renvoyant le \texttt{BacktestResult}, les ordres et les exécutions.
\bfield{Règles} Les gestionnaires ne programment que des événements futurs~; les ordres ne sont exécutés qu'après leur accusé de réception~; toute hypothèse d'exécution réside dans un \omterm{def:rs:event-driven-backtests:fill}{modèle d'exécution}~; la
politique de présence dans la barre est énoncée~; l'exécution est déterministe.
\bfield{Tests d'acceptation} \texttt{code/firm/evbt/tests/}~: \omterm{def:rs:event-driven-backtests:fill}{exécutions au contact}, par pénétration et plafonnées sur des barres témoins avec leurs cycles de vie~; un ordre au marché à
l'ouverture suivante~; une latence qui manque une barre~; des annulations~; la politique optimiste exécutant un ordre présent pendant un instant~; les frais, le P\&L et un split~; deux exécutions identiques.
\bfield{Pour aller plus loin} Dividendes, demi-journées, enchères et suspensions en événements de calendrier~; ordres stop~; plusieurs instruments sur une même horloge.
\end{build}

\begin{omsources}
\item Zipline (backtester Python open source), \texttt{zipline/finance/slippage.py}~: \texttt{VolumeShareSlippage}, limite de volume par défaut de 2.5~\% de chaque barre.
\end{omsources}

\section{Exercices}

\begin{exercise}[$\star$]\label{exo:rs:event-driven-backtests:1}
Un achat à cours limité à 99.00 est en attente. La barre suivante a un cours d'ouverture de 99.20, un plus haut de 99.30, un plus bas de 99.00 et un cours de clôture de 99.10. Est-il exécuté au contact~? Et sous une \omterm{def:rs:event-driven-backtests:fill}{exécution par pénétration}
d'un tick, le tick valant 0.01~?
\end{exercise}

\begin{exercise}[$\star$]\label{exo:rs:event-driven-backtests:2}
Un ordre portant sur 5\,000 actions fait face à des barres de 40\,000, 10\,000 et 60\,000 actions avec un plafond de volume de 2.5~\%. Quelle quantité est exécutée dans chaque barre~?
\end{exercise}

\begin{exercise}[$\star$]\label{exo:rs:event-driven-backtests:3}
Dessiner le cycle de vie d'un ordre portant sur 300 actions, exécuté à hauteur de 100, puis annulé avec 200 restantes.
\end{exercise}

\begin{exercise}[$\star\star$]\label{exo:rs:event-driven-backtests:4}
Pourquoi les \omterm{def:rs:event-driven-backtests:fill}{exécutions par pénétration} suivies de mouvements défavorables du cours sont-elles plus nombreuses que les \omterm{def:rs:event-driven-backtests:fill}{exécutions au contact}~? Relier cela à la file au cours de l'ordre.
\end{exercise}

\begin{exercise}[$\star\star$]\label{exo:rs:event-driven-backtests:5}
Avec une latence de 5 secondes et des barres d'une minute, pourquoi le teneur de cotation n'est-il jamais exécuté sous la politique prudente~? Comment se comporterait une stratégie qui n'annule que lorsque
le cours bouge~?
\end{exercise}

\begin{exercise}[$\star\star$]\label{exo:rs:event-driven-backtests:6}
Un split de deux pour un prend effet alors qu'un ordre de vente à cours limité à 120~\$ portant sur 100 actions est en attente et que la position est de 100 actions. Que doit faire le moteur, et que
se passe-t-il s'il ne le fait pas~?
\end{exercise}

\begin{exercise}[$\star\star\star$]\label{exo:rs:event-driven-backtests:7}
\emph{Programmation.} Lancer le teneur de cotation avec des ordres au cours de clôture (et non à un tick de distance) sous les modèles au contact et par pénétration. Que deviennent les taux d'exécution et
les mark-outs~?
\end{exercise}

\begin{exercise}[$\star\star\star$]\label{exo:rs:event-driven-backtests:8}
\emph{Trouver la faille.} «~Notre \omterm{def:rs:vectorised-backtests:backtest}{backtest} exécute les ordres à cours limité dès que le plus bas de la barre les atteint, et la stratégie a gagné de l'argent chaque semaine pendant deux ans.~»
\end{exercise}

\section{Problème~: l'ordre à cours limité qui était toujours exécuté}

\begin{problem}[{Problème du week-end --- une stratégie passive, trois modèles d'exécution}]\label{pb:rs:event-driven-backtests:1}
Le teneur de cotation du chapitre sur quatre journées simulées des barres d'une minute de \texttt{firm.tape}.

\medskip\noindent\textbf{Partie I --- Le moteur.}
\begin{enumerate}
\item Dans quel ordre \texttt{firm.evbt} traite-t-il la clôture d'une barre~: exécutions, valorisations, stratégie~? Pourquoi cet ordre~?
\item Quand un ordre soumis peut-il pour la première fois être exécuté~?
\item Qu'est-ce qui rend deux exécutions identiques~?
\item Que renvoie le moteur~?
\end{enumerate}

\medskip\noindent\textbf{Partie II --- \omterm{def:rs:event-driven-backtests:fill}{Modèles d'exécution}.}
\begin{enumerate}[resume]
\item Donner les lots exécutés, le taux d'exécution, le P\&L et le mark-out d'une journée sous chaque \omterm{def:rs:event-driven-backtests:fill}{modèle d'exécution} sans latence.
\item Quelle hypothèse sur la file chaque modèle fait-il~?
\item Pourquoi le P\&L de l'\omterm{def:rs:event-driven-backtests:fill}{exécution par pénétration} est-il négatif chaque jour alors que celui de l'\omterm{def:rs:event-driven-backtests:fill}{exécution au contact} ne l'est pas~?
\item Quelle est la version au niveau 1 de l'idée, et que rapporte-t-elle~?
\end{enumerate}

\medskip\noindent\textbf{Partie III --- Latence.}
\begin{enumerate}[resume]
\item Que se passe-t-il avec une latence de 5 secondes sous la politique prudente~?
\item Et sous la politique optimiste, pour chaque \omterm{def:rs:event-driven-backtests:fill}{modèle d'exécution}~?
\item Que vous apprend la différence entre les deux politiques sur la stratégie~?
\item Comment modifieriez-vous la stratégie pour que la latence compte moins~?
\end{enumerate}

\medskip\noindent\textbf{Partie IV --- Le verdict.}
\begin{enumerate}[resume]
\item Énoncer le \emph{résultat nommé}~: le taux d'exécution et le P\&L de la stratégie passive par barres sous \omterm{def:rs:event-driven-backtests:fill}{exécution au contact}, par pénétration et plafonnée, avec et sans
latence.
\item Quel résultat présenteriez-vous à un gérant de portefeuille, et avec quelle réserve~?
\item Qu'apporterait un rejeu de niveau 3 (chapitre~18)~?
\item Que signifie le plafond par défaut de Zipline pour une stratégie qui échange 5~\% du volume~?
\item Quelles opérations sur titres doivent figurer dans la file d'événements pour cette stratégie~?
\item Pourquoi un taux d'exécution de 70~\% pour des ordres situés à un tick de distance est-il suspect~?
\item Que journaliseriez-vous lors d'une exécution en direct pour calibrer le \omterm{def:rs:event-driven-backtests:fill}{modèle d'exécution}~?
\item En une phrase~: qu'est-ce que le \omterm{def:rs:event-driven-backtests:fill}{modèle d'exécution}~?
\end{enumerate}
\end{problem}

\section{Questions d'entretien}

\begin{interviewq}[$\star$ \iqroles{développeur, chercheur}]\label{iq:rs:event-driven-backtests:1}
Décrivez l'architecture d'un \omterm{def:rs:event-driven-backtests:evbt}{backtest événementiel}.
\end{interviewq}

\begin{interviewq}[$\star\star$ \iqroles{chercheur, trader}]\label{iq:rs:event-driven-backtests:2}
Comment simuleriez-vous les exécutions d'ordres à cours limité sur des données en barres, et quels biais chaque approche comporte-t-elle~?
\end{interviewq}

\begin{interviewq}[$\star\star$ \iqroles{développeur}]\label{iq:rs:event-driven-backtests:3}
Comment rendez-vous un \omterm{def:rs:vectorised-backtests:backtest}{backtest} déterministe~?
\end{interviewq}

\begin{interviewq}[$\star\star$ \iqroles{chercheur}]\label{iq:rs:event-driven-backtests:4}
Qu'est-ce que la sélection adverse pour un ordre passif, et comment la mesureriez-vous dans un \omterm{def:rs:vectorised-backtests:backtest}{backtest}~?
\end{interviewq}

\begin{interviewq}[$\star\star$ \iqroles{développeur, trader}]\label{iq:rs:event-driven-backtests:5}
Comment un \omterm{def:rs:vectorised-backtests:backtest}{backtest} doit-il traiter une division d'actions pendant que des ordres sont en cours~?
\end{interviewq}

\begin{interviewq}[$\star\star\star$ \iqroles{chercheur}]\label{iq:rs:event-driven-backtests:6}
Quand un \omterm{def:rs:vectorised-backtests:backtest}{backtest} de niveau 2 ne suffit-il pas~? Donnez un test qui le montre.
\end{interviewq}
